\section{Postać Jordana macierzy.}

W tej sekcji zakładamy,
	że ustalone mamy ciało $\mathbb{K}$,
	które jest równe $\mathbb{R}o$ lub $\mathbb{C}$.

\begin{definition}
	\textit{Wektorem własnym} macierzy $A \in \mathbb{K}^{n \times n}$
		nazywamy niezerowy wektor $v \in \mathbb{K}^{n \times 1}$,
		że dla pewnego $\lambda \in \mathbb{K}$
	\begin{equation}
		Av = \lambda v.
	\end{equation}
	$\lambda$ nazywamy \textit{wartością własną} odpowiadającą wektorowi własnemu $v$.
\end{definition}

Załóżmy,
	że daną mamy macierz $A in \mathbb{K}^{n \times n}$
	i chcemy wyznaczyć jej wektory i wartości własne.
Jak możemy to zrobić?
Przekształćmy równanie definiujące te pojęcia.
\begin{align}
	Av = \lambda v, \\
	\textbf{0} = \lambda v - Av, \\
	\textbf{0} = \lambda I v - Av, \\
	\textbf{0} = (\lambda I - A) v.
\end{align}
Oznacza to,
	że $v$ jest wektorem własnym macierzy $A$ z odpowiadającą mu wartością własną $\lambda$
	wtedy i tylko wtedy gdy $(\lambda I - A) v = \textbf{0}$,
	czyli wtedy i tylko wtedy gdy $v \in \ker(\lambda I - A)$.

Skoro interesują nas tylko takie $v$,
	które są niezerowe,
	to warunek ten może być spełniony wtedy i tylko wtedy,
	gdy jądro $\ker(\lambda I - A)$ jest nietrywialne.
Zachodzi to dokładnie wtedy,
	gdy $A - \lambda I$ jest macierzą nieodwracalną,
	czyli
\begin{equation}\label{eq:eigen-condition}
	\det(\ker(\lambda I - A)) = 0.
\end{equation}

Aby wyznaczyć wartości własne $A$ musimy więc sprawdzić,
	dla jakich $\lambda \in \mathbb{K}$ spełnione jest równanie \ref{eq:eigen-condition}.
Okazuje się to dość proste, gdyż zawsze jest to równanie wielomianowe (zadanie \ref{ex:characteristic-eq-is-poly})
Kiedy to ustalimy,
	wyznaczyć możemy $\ker(A - \lambda I)$.
Każdy wektor w tej przestrzeni liniowej będzie wektorem własnym $A$.
Wielomian,
	który pojawia się w tym równaniu ma swoją specjalną nazwę.

\begin{definition}
	\textit{Wielomianem charakterystycznym} macierzy $A \in \mathbb{K}^{n \times n}$
		nazywamy wielomian
	\begin{equation}
		\chi_A(t) = \det(t I - A).
	\end{equation}
\end{definition}

\begin{definition}
	Krotnością algebraiczną wartości własnej $\lambda$ macierzy $A \in \mathbb{K}^{n \times n}$
		nazywamy krotność pierwiastka $\lambda$ wielomianu $\chi_A(\lambda)$.
\end{definition}

\begin{example}
	Niech
	\begin{equation}
		A = \left[\begin{matrix}
		\end{matrix}\right].
	\end{equation}
	Wyznaczymy wartości własne macierzy $A$ oraz ich krotności.
\end{example}

% TODO: 0 vs \textbf{0}


\begin{definition}
	Macierz $A \in \mathbb{K}^{n \times n}$ nazywamy \textit{diagonizowalna},
		jeśli istnieje macierz odwracalna $X \in \mathbb{K}^{n \times n}$
		i macierz diagonalna $J \in \mathbb{K}^{n \times n}$,
		że
	\begin{equation}
		A = X J X^{-1}.
	\end{equation}
\end{definition}

Załóżmy,
	że mamy daną macierz $A$.


\begin{example}
	Macierz
	\begin{equation}
		A = \left[\begin{matrix}
			1 & 1 & 0 \\
			1 & 1 & 0 \\
			0 & 0 & 2
		\end{matrix}\right]
	\end{equation}
	jest diagonizowalna.
	Wyznaczmy jej wartości własne.
\end{example}

\begin{example}
	Macierz
	\begin{equation}
		A = \left[\begin{matrix}
			2 & 1 & 0 \\
			0 & 2 & 1 \\
			0 & 0 & 2
		\end{matrix}\right]
	\end{equation}
	nie jest diagonalizowalna.
	
\end{example}

\begin{definition}
	Dla dowolnego $m \in \mathbb{Z}_+$ i dowolnego $\lambda \in \mathbb{K}$
		\textit{klatką Jordana} nazywamy macierz $J_m(\lambda) \in \mathbb{K}^{m \times m}$
	postaci
	\begin{equation}
		J_m(\lambda) = \begin{matrix}
			\lambda & 1 & 0 & \dots & 0 & 0
		\end{matrix}
	\end{equation}
\end{definition}

\begin{enumerate}
	\item oblicz wartości własne $A: \lambda_1, \dots, \lambda_k$,
	\item oblicz krotości algebraiczne $\lambda_1, \dots, \lambda_k$
		(krotości zer w wielomianie charakterystycznym $A$),
	\item oblicz krotności geometryczne $\lambda_1, \dots, \lambda_k$
		(krotność geometryczna $\lambda_i$ to $\dim(\ker(A - \lambda_i I)))$),
	\item (opcjonalnie) oblicz $\dim((A - \lambda_i I)^2), \dim((A - \lambda_i I)^3), \dim((A - \lambda_i I)^4), \dots$, aż wymiar przestanie się zmieniać.
\end{enumerate}

\section{Zadania.}
